%%%%%%%%%%%%%%%%%%%%%%%%%%%%%%%%%%%%%%%%%%%%%%%%%%%%%%%%%%%%%%%%%%%%
%% I, the copyright holder of this work, release this work into the
%% public domain. This applies worldwide. In some countries this may
%% not be legally possible; if so: I grant anyone the right to use
%% this work for any purpose, without any conditions, unless such
%% conditions are required by law.
%%%%%%%%%%%%%%%%%%%%%%%%%%%%%%%%%%%%%%%%%%%%%%%%%%%%%%%%%%%%%%%%%%%%

\documentclass{beamer}
\usetheme[faculty=ped,basePath=../fibeamer,logo=../fibeamer/logo/mu/Logo-UC-3.png]{fibeamer}
\graphicspath{{../resources/}}

\usepackage{algorithm}
\usepackage{algpseudocode}

\usepackage[utf8]{inputenc}
\usepackage[
  main=english, %% By using `czech` or `slovak` as the main locale
                %% instead of `english`, you can typeset the
                %% presentation in either Czech or Slovak,
                %% respectively.
  czech, slovak %% The additional keys allow foreign texts to be
]{babel}        %% typeset as follows:
%%
%%   \begin{otherlanguage}{czech}   ... \end{otherlanguage}
%%   \begin{otherlanguage}{slovak}  ... \end{otherlanguage}
%%
%% These macros specify information about the presentation
\title{Ciclos anidados} %% that will be typeset on the
\subtitle{Semestre II-2026} %% title page.
\author{Profesora: Andreina Cota Vidaurre}
%% These additional packages are used within the document:
\usepackage{ragged2e}  % `\justifying` text
\usepackage{booktabs}  % Tables
\usepackage{tabularx}
\usepackage{tikz}      % Diagrams
\usetikzlibrary{calc, shapes, backgrounds}
\usetikzlibrary{positioning, arrows.meta, shadows.blur, shapes.geometric}
\usepackage{amsmath, amssymb}
\usepackage{url}       % `\url`s
\usepackage{listings}  % Code listings
\usepackage{graphicx}
\usepackage{amsmath}
\usepackage[T1]{fontenc}
\usepackage{xcolor}
\usepackage{minted}
% \usepackage{adjustbox}


% Definicion de pseudocodigo en espaniol
\lstdefinelanguage{PseudocodigoES}{
    morekeywords={INICIO,FIN,PARA,HACER,SI,ENTONCES,FIN_SI,FIN_PARA,IMPRIMIR},
    sensitive=true,
    morecomment=[l]{//},
    morestring=[b]"
}

% Definicion de pseudocodigo en ingles
\lstdefinelanguage{PseudocodeEN}{
    morekeywords={BEGIN,END,FOR,DO,IF,THEN,END_IF,END_FOR,PRINT},
    sensitive=true,
    morecomment=[l]{//},
    morestring=[b]"
}

\lstset{
    basicstyle=\ttfamily\small,
    keywordstyle=\bfseries,
    columns=fullflexible,
    keepspaces=true,
    showstringspaces=false,
    frame=single,
    xleftmargin=0em,
    framexleftmargin=0em,
    framesep=2pt,
    gobble=4,
    resetmargins=true,
    aboveskip=0pt,
    belowskip=0pt
}

\lstdefinestyle{pythonstyle}{
    language=Python,
    basicstyle=\ttfamily\small,
    keywordstyle=\bfseries,
    commentstyle=\itshape,
    stringstyle=\ttfamily,
    showstringspaces=false,
    columns=fullflexible,
    keepspaces=true,
    frame=single,
    breaklines=true,
    xleftmargin=0em,
    framexleftmargin=0em,
    framesep=2pt,
    aboveskip=0pt,
    belowskip=0pt,
    gobble=4,
    resetmargins=true,
    emph={print,input,int,float,str,bool,len,range},
    emphstyle=\bfseries
}

\lstset{style=pythonstyle}

% Estilos del diagrama
\tikzstyle{iniciofin} = [
    ellipse,
    minimum width=2.5cm,
    minimum height=1cm,
    text centered,
    draw=black,
    fill=blue!25
]

\tikzstyle{proceso} = [
    rectangle,
    rounded corners,
    minimum width=3cm,
    minimum height=1cm,
    text centered,
    draw=black,
    fill=cyan!20
]

\tikzstyle{decision} = [
    diamond,
    aspect=2,
    text centered,
    draw=black,
    fill=yellow!35
]

\tikzstyle{flecha} = [
    thick,
    ->,
    >=Stealth
]

% \definecolor{green_python}{HTML}{52A736}

% \newcolumntype{Y}{>{\raggedright\arraybackslash}X}

\frenchspacing
\begin{document}
  \shorthandoff{-}
  \frame[c]{\maketitle}

% \begin{darkframes}

    \begin{frame}[fragile, t]{\mintinline{python}{for + if}}
        \begin{itemize}
            \item \mintinline{python}{for + if}: Recorre el elemento por elemento y, en cada iteración, verifica si se cumple la condición. Si se cumple, ejecuta el bloque interno.
            \begin{minted}[xleftmargin=-2cm, frame=leftline, fontsize=\small]{python}
            for item in secuencia:
                if condicion:
                    instrucciones_if
                instrucciones_for
            \end{minted}
        \end{itemize}
    \end{frame}

    \begin{frame}[fragile, t]{\mintinline{python}{while + if}}
        \begin{itemize}
            \item \mintinline{python}{while + if}: Repite el bloque mientras la condición principal sea verdadera y, en cada vuelta, verifica una condición adicional. Son dos condiciones independientes: una controla cuándo ocurre el ciclo, la otra decide qué hace en cada vuelta.
            \begin{minted}[xleftmargin=-2cm, frame=leftline, fontsize=\small]{python}
            while condicion_while:
                if condicion_if:
                    instrucciones_if
                instrucciones_while
            \end{minted}
        \end{itemize}
    \end{frame}

    \begin{frame}[fragile, t]{Ciclos anidados: \mintinline{python}{for + for}}
        \begin{itemize}
            \item \mintinline{python}{Ciclos anidados}: Un ciclo dentro de otro
            \begin{minted}[xleftmargin=-2cm, frame=leftline, fontsize=\small]{python}
            for i in range(n):
                # Ciclo interno que se ejecuta n veces
                for j in range(m):
                    # Ciclo interno que se ejecuta m veces
                    instrucciones_for_interno
                # Fin del ciclo interno
                instrucciones_for_externo
            # Fin del ciclo externo
            \end{minted}
            \textcolor{blue}{El ciclo interno se ejecuta completo en cada iteración del ciclo externo}
        \end{itemize}
    \end{frame}

    \begin{frame}[fragile, t]{Ciclos anidados: \mintinline{python}{while + while}}
        \begin{itemize}
            \item \mintinline{python}{Ciclos anidados}: Un ciclo dentro de otro
            \begin{minted}[xleftmargin=-2cm, frame=leftline, fontsize=\small]{python}
            while condicion_while_externo:
                # Ciclo externo que se ejecuta mientras la condición 
                # sea verdadera
                while condicion_while_interno:
                    # Ciclo interno que se ejecuta mientras la 
                    # condición sea verdadera
                    instrucciones_while_interno
                # Fin del ciclo interno
                instrucciones_while_externo
            # Fin del ciclo externo
            \end{minted}
            % \textcolor{blue}{El ciclo interno se ejecuta mientras la condición sea verdadera en cada iteración del ciclo externo}
        \end{itemize}
    \end{frame}

    \begin{frame}[fragile, t]{Ciclos anidados: \mintinline{python}{for + while}}
        \begin{itemize}
            \item \mintinline{python}{Ciclos anidados}: Un ciclo dentro de otro
            \begin{minted}[xleftmargin=-2cm, frame=leftline, fontsize=\small]{python}
            for i in range(n):
                # Ciclo interno que se ejecuta n veces
                while condicion_while_interno:
                    # Ciclo interno que se ejecuta mientras la 
                    # condición sea verdadera
                    instrucciones_while_interno
            \end{minted}
            % \textcolor{blue}{El ciclo interno se ejecuta mientras la condición sea verdadera en cada iteración del ciclo externo}
        \end{itemize}
    \end{frame}

    \begin{frame}[fragile, t]{Ciclos anidados: \mintinline{python}{while + for}}
        \begin{itemize}
            \item \mintinline{python}{Ciclos anidados}: Un ciclo dentro de otro
            \begin{minted}[xleftmargin=-2cm, frame=leftline, fontsize=\small]{python}
            while condicion_while:
                # Ciclo externo que se ejecuta mientras la condición 
                # sea verdadera
                for i in range(n):
                    # Ciclo interno que se ejecuta n veces
                    instrucciones_for_interno
            \end{minted}
            % \textcolor{blue}{El ciclo interno se ejecuta completo en cada iteración del ciclo externo}
        \end{itemize}
    \end{frame}

    \begin{frame}[fragile, t]{Completar el código \mintinline{python}{for + if}}
        En cada ejercicio encontrarás código incompleto. Tu tarea es completar los espacios en blanco para que el programa funcione correctamente.
        \begin{itemize}
            \item Completa el código para imprimir una cuenta regresiva del 10 al 1. Cuando llegue a 5, además del número, imprime '¡Mitad del camino!'. Al terminar, imprime '¡Despegue!'
            \begin{minted}[xleftmargin=-2cm, frame=leftline, fontsize=\small]{python}
            for i in range(________, ________, ________):
                print(i)
                if i ________ 5:
                    print('¡Mitad del camino!')
            
            print('¡Despegue!')
            \end{minted}
            \textcolor{blue}{Salida esperada}: 10 9 8 7 6 5 !Mitad del camino! 4 3 2 1 !Despegue!
        \end{itemize}
    \end{frame}

    \begin{frame}[fragile,t]{Completar el código \mintinline{python}{while + if}}
        \begin{itemize}
            \item Un vehículo blindado parte con 100 litros de combustible y consume 8 litros por kilómetro. Completa el código para simular el avance hasta que se quede sin combustible. Si el combustible baja por debajo de 25 litros, imprime una alerta.
            \begin{minted}[xleftmargin=-2cm, frame=leftline, fontsize=\small]{python}
            combustible = 100
            km = 0
            while combustible ________ 0:
                km ________ 1
                combustible ________ 8
                if combustible ________ 25:
                    print(f'km {km}: Alerta — {combustible}L restantes')
                else:
                    print(f'km {km}: {combustible}L')
            print(f'El vehículo avanzó {km} km')
            \end{minted}
            \textcolor{blue}{SE:} km 1: 92L .. km 10: 20L (Alerta)  ... |  El vehículo avanzó 13 km
        \end{itemize}
    \end{frame}

    \begin{frame}[fragile, t]{Completar el código \mintinline{python}{for + for}: Ciclos anidados}
        \begin{itemize}
            \item Completa el código para imprimir cuántos ejercicios realiza un estudiante en cada ronda. Las repeticiones son: número de ronda × número de ejercicio.
            \begin{minted}[xleftmargin=-2cm, frame=leftline, fontsize=\small]{python}
            for ronda in range(1, 4):         # 3 rondas
                for ejercicio in range(1, 5): # 4 ejs por ronda
                    reps = ________ * ________
                    print(f'Ronda {ronda} - \ 
                            Ejercicio {ejercicio}: {reps} reps')
            \end{minted}
            \textcolor{blue}{Salida esperada}: Ronda 1 - Ejercicio 1: 1 reps  ...  Ronda 1 - Ejercicio 4: 4 reps  |  Ronda 2 - Ejercicio 1: 2 reps  ...
        \end{itemize}
    \end{frame}

    \begin{frame}[fragile, t]{Completar el código \mintinline{python}{while + while}: Ciclos anidados}
        \begin{itemize}
            \item Completa el código para recorrer una matriz de 3x4.
            \begin{minted}[xleftmargin=-2cm, frame=leftline, fontsize=\small]{python}
            i = 0
            while i < _______:         # 3 rondas
                j = 0
                while j < _______: # 4 ejs por ronda
                    print(f'Elemento ({i}, {j})')
                    j ________ 1
                i ________ 1
            \end{minted}
            \textcolor{blue}{Salida esperada}: Elemento (0, 0)  ...  Elemento (2, 2)
        \end{itemize}
    \end{frame}

    \begin{frame}[fragile, t]{Ciclos anidados}
        \begin{itemize}
            \item Un \textbf{ciclo anidado} corresponde a una una estructura iterativa contenida dentro de otra estructura iterativa.
            \item El \textbf{ciclo externo} ontrola el recorrido de la estructura principal, mientras que el \textbf{ciclo interno} procesa sus elementos internos.
            \item Por cada iteracion del ciclo externo, el ciclo interno realiza nuevamente su recorrido.
            \item Los ciclos anidados son especialmente útiles para recorrer y procesar estructuras de datos anidadas, como listas de listas, listas de tuplas o matrices.
        \end{itemize}
    \end{frame}

    \begin{frame}[fragile, t]{Ciclos anidados}
        \begin{itemize}
            \item Es posible combinar distintos tipos de ciclos: \mintinline{python}{for + for}, \mintinline{python}{while + while}, \mintinline{python}{for + while}, \mintinline{python}{while + for}, etc.
            \item \mintinline{python}{for} es apropiado cuando se conoce la estructura o cantidad de elementos que se desea recorrer.
            \item \mintinline{python}{while} es apropiado cuando la repeticion depende del cumplimiento de una condicion.
            \item En ciclos \mintinline{python}{while} anidados, cada nivel requiere controlar e inicializar correctamente sus variables para evitar ciclos infinitos.
            \item Independientemente del tipo de ciclo utilizado, el principio de anidamiento es el mismo: \textbf{recorrer diferentes niveles de una estructura o problema de manera organizada}.
        \end{itemize}
    \end{frame}


\end{document}
